\section{Introduction}
\label{sec:introduction}

% ---------------------------------------------------------------
%  PARAGRAPH 1: Context — Paper 1 summary
% ---------------------------------------------------------------

In a companion paper~\cite{KharaziLeytonOrtega2026a} we introduced
stereographic quantum signal processing (stereo-QSP), showing that
the boundedness constraint of QSP~\cite{LowChuang2017,GSLW2019}
can be lifted by changing the classical encoding and decoding
interface rather than the quantum algorithm.  A single-qubit signal
$r \in [0,\infty)$ is compressed into the bounded QSP domain via
$\rtilde = r/\sqrt{1+r^2}$, and Pauli-measurement decoding
recovers the ratio $P(\rtilde)/Q(\rtilde)$---an unbounded rational
function of the original parameter.  The key structural result is
that the stereographic and standard QSP signal operators are
unitarily equivalent via $V = \mathrm{diag}(1,i)$, so all existing
phase-finding algorithms and convergence results apply directly.

% ---------------------------------------------------------------
%  PARAGRAPH 2: The multi-qubit question
% ---------------------------------------------------------------

The natural next question---and the subject of this paper---is
whether stereographic encoding extends to multi-qubit Hamiltonians.
In the standard quantum singular value transformation
(QSVT)~\cite{GSLW2019} framework, a Hamiltonian
$H$ with $\|H\| = \alpha$ is embedded into a larger unitary via an
$(\alpha, m, \epsilon)$-block encoding, which stores $H/\alpha$ in
the top-left block.  The normalization factor
$\alpha \geq \|H\|$ inflates the circuit depth by a factor
of~$\alpha$ and forces the QSP polynomial to approximate
$f(\alpha\,\cdot\,)$ on $[-1,1]$.  Stereographic block encoding
replaces this with an exact, single-ancilla embedding that places
each eigenvalue directly onto the Bloch sphere, removing the
normalization entirely.

% ---------------------------------------------------------------
%  PARAGRAPH 3: The eigenbasis requirement — honest assessment
% ---------------------------------------------------------------

This advantage comes at a price: the construction requires
knowledge of the eigenbasis.  Whereas standard block encoding
embeds $H/\alpha$ into a unitary without distinguishing eigenspaces,
the stereographic version applies a different encoding rotation in
each eigenspace, conditioning on the eigenstates of~$H$.  For
generic Hamiltonians, the eigenbasis is the object one wishes to
compute, and the construction becomes circular.  Stereographic block
encoding is therefore restricted to \emph{structured Hamiltonians}
whose eigenbasis is known or efficiently computable: diagonal
operators, integrable models, or operators whose spectral
decomposition is available as a quantum circuit.  A central goal of
this paper is to delineate precisely which Hamiltonian classes are
tractable and what the construction costs in each case.

% ---------------------------------------------------------------
%  PARAGRAPH 4: Contributions of this paper
% ---------------------------------------------------------------

Our contributions are as follows.  We give a formal treatment of
stereographic block encoding, proving that it is exact (no
$\epsilon$ parameter), uses a single ancilla qubit, and produces
circuit depth $\mathcal{O}(d)$ independent of $\|H\|$
(Section~\ref{sec:block-encoding}).  We show explicitly that the
single-qubit QSP protocol lifts to a multi-qubit eigenvalue
transformation, with phase gates commuting through the system
register and the signal operator decomposing as a direct sum over
eigenspaces (Section~\ref{sec:qsvt}).  We construct explicit
circuits for three Hamiltonian classes of increasing
complexity---diagonal operators, Pauli-$Z$ sums with reversible
adders, and the Heisenberg model with a Givens-rotation change of
basis---providing gate counts, depth analysis, and comparison with
standard block encoding (Section~\ref{sec:circuits}).  We analyze
the decoding protocol when the system register is in a superposition
of eigenspaces, characterizing the post-measurement state and the
measurement overhead (Section~\ref{sec:decoding}).  Finally, we
present expanded numerical simulations---Qulacs-verified circuits
for the Heisenberg model at multiple parameter values, gate-count
benchmarks, and noise analysis---establishing the practical scope of
the approach (Section~\ref{sec:numerics}).

% ---------------------------------------------------------------
%  PARAGRAPH 5: Paper organization
% ---------------------------------------------------------------

The remainder of this paper is organized as follows.
Section~\ref{sec:review} provides a self-contained review of
stereographic QSP from the companion paper.
Section~\ref{sec:block-encoding} defines stereographic block
encoding and compares it with standard block encoding.
Section~\ref{sec:qsvt} develops the full QSVT protocol.
Section~\ref{sec:circuits} constructs explicit circuits for three
Hamiltonian classes.
Section~\ref{sec:decoding} analyzes the decoding and measurement
protocol.
Section~\ref{sec:numerics} presents numerical results.
Section~\ref{sec:discussion} discusses connections to prior work,
tractable Hamiltonian classes, and open problems.
